
(4.8.1), on obtient :
@@M0@@
soit d'après 3.6 :
@@M1@@

Pour montrer la relation ii), il faut montrer que :
@@M2@@
Posons @@M3@@. On a :
@@M4@@
Or la relation i) donne :
@@M5@@
Compte tenu de (4.8.1), on obtient donc :
@@M6@@
d'où :
@@M7@@

Soit @@M8@@ l'image de @@M9@@ dans @@M10@@, et @@M11@@ celle de @@M12@@ dans @@M13@@. On a donc : @@M14@@. On peut alors appliquer (VI 5.11) (qui reste valable après tensorisation par @@M15@@, car sa démonstration s'appuie seulement sur la structure de @@M16@@ comme @@M17@@-algèbre libre), et on peut donc trouver @@M18@@ tel que :
@@M19@@
En appliquant @@M20@@, on trouve, compte tenu de ce que @@M21@@ :
@@M22@@
d'où, puisque @@M23@@ :
@@M24@@
ce qui achève la démonstration.

